% !TeX root = ../深度学习与大语言模型技术原理.tex
% 由主文档通过 \input 载入；请编译主文档。
\setcounter{chapter}{15} % 保留原书章号。
\chapter{工具接口、MCP 与结构化生成}
\label{ch:tools}

\section{从输出文本到可执行动作}

工具调用（tool calling）或函数调用（function calling）使模型输出工具名和参数，由外部程序执行。模型生成的调用只是动作建议，执行程序才负责参数解析、权限校验、资源分配和真实副作用。工具模式（tool schema）应明确输入类型、必填字段、枚举、返回结构及错误类别；冗长描述无法替代明确类型约束。

例如检索工具的输入可以包含查询字符串、最大结果数和时间范围，输出包含文档标识、片段、来源与时间。调用记录应保存请求标识、工具版本、参数、开始结束时间、结果及状态。把普通错误描述成正常结果会误导下一步推理，因此应区分成功、可重试故障、不可重试故障和状态未知。

\emph{幂等性（idempotency）}要求同一操作重复执行不会产生额外的不可忽略副作用。读取通常容易重试，转账、发送消息或提交订单则需要请求键、去重和状态查询。超时只说明调用方未及时收到结果，并不证明服务端没有执行；状态未知时盲目重试可能重复写入。长任务还应支持取消、结果恢复和明确的部分成功语义。

\section{MCP 的宿主、客户端与服务端}

模型上下文协议（Model Context Protocol, MCP）标准化应用与外部工具、资源之间的连接。按本书引用的2025-06-18版规范，宿主（host）协调模型、用户与多个客户端（client），每个客户端连接相应服务端（server）；服务端暴露工具、资源和提示模板等能力\cite{mcp}。协议解决发现、消息交换和能力协商，不替代模型推理算法或应用的权限判断。

工具用于请求计算或动作，资源（resources）用于提供可读取内容，提示模板（prompts）封装可选择的交互模板。连接初始化时协商版本与能力，之后按约定交换消息。该版本包含标准输入输出与 Streamable HTTP 等传输方式；传输、鉴权与会话管理的细节应按部署采用的规范版本核对，不能混用不同版本示例。

宿主必须决定哪些能力暴露给模型、哪些资源可访问，以及何时实际执行动作。服务端返回的文档或描述可能包含不可信内容，即使连接成功也不代表其文本具有指令权限。凭证应由受控程序管理，工具不应因为模型在参数中声称“已获授权”就扩大权限。

\section{JSON 模式与生成后校验}

结构化输出（structured output）规定生成结果满足某种可解析形式，例如 JSON 对象或受限代码。JSON 语法合法只说明括号、字符串和分隔符正确；JSON Schema 还可规定字段类型、必填项、枚举和数值范围。字段“日期”是合法字符串，不保证它是真实日期；金额在范围内，也不保证操作被授权。

仅用提示词要求格式，模型仍可能生成前后说明或遗漏字段。生成后校验能发现不合格结果，再进行有限次修复或返回明确错误；修复过程本身也必须受预算和权限控制。对关键动作，业务规则应在确定性程序中再次验证，并把解析失败、语义失败和权限失败分开统计。

\section{语法约束解码与可行前缀}

约束解码（constrained decoding）在生成每个词元前，维护解析状态 $s$，计算允许集合 $\mathcal A(s)$，然后取
\begin{equation}
\widetilde p(v\mid s)=\frac{p(v\mid s)\mathbf1\{v\in\mathcal A(s)\}}{
\sum_{u\in\mathcal A(s)}p(u\mid s)}.
\label{eq:constrained}
\end{equation}
它要求分母非零，并应只保留能够继续扩展为合法完整结果的前缀。词元可能包含多个字符或部分字节，不能把字符级合法性直接等同于词元级合法性。

有限状态机（finite-state machine, FSM）适合正则形式；任意深度嵌套通常需要栈、上下文无关语法解析器或限定深度展开。结束词元只能在合法完成状态允许，长度预算耗尽时也不能把未闭合结果报告为成功。不同后端支持的模式子集可能不同，应把模式本身纳入测试。

逐步重归一化一般会改变概率分布；它不自动等于原模型在“最终结果满足全局约束”条件下的精确条件分布，因为后者还包含各前缀未来完成的概率。约束解码提供的是相应形式上的保证，不是内容正确性或全局最优性保证。

